\documentclass{article}
\usepackage[T1]{fontenc}
\usepackage{lmodern}
\usepackage{graphicx}
\usepackage{amsmath,amssymb}
\usepackage[a4paper,margin=2cm]{geometry}
\usepackage[absolute,overlay]{textpos}
\setlength{\TPHorizModule}{1mm}
\setlength{\TPVertModule}{1mm}
\usepackage[indonesian]{babel}
\usepackage{hyperref}
\setlength{\emergencystretch}{2em}
% unit: Halaman 53

\begin{document}

% === Halaman 53 (python/native) ===

53

1. Confusion

• Prinsip ini menyembunyikan hubungan statistik yang ada diantara 

plainteks, cipherteks, dan kunci. 


• Prinsip confusion membuat kriptanalis frustasi untuk mencari pola-

pola statistik yang muncul di dalam cipherteks. 

• Confusion dapat direalisasikan dengan menggunakan teknik substitusi 

yang nirlanjar (non-linear), biasanya dengan operasi lookup table.

\end{document}
